\section{Các nghiên cứu liên quan}
\label{sec:related-work}

\subsection{Kiểm soát truy cập dựa trên vai trò}
\label{subsec:rbac}

RBAC là một mô hình kiểm soát truy cập được sử dụng rộng rãi, trong đó quyền
được gán cho vai trò thay vì gán trực tiếp cho người dùng. Mô hình RBAC ban đầu
do Ferraiolo và Kuhn đề xuất \citep{Ferraiolo1992RBAC} cùng họ mô hình RBAC96
của Sandhu và cộng sự \citep{Sandhu1996RBAC} cung cấp nền tảng khái niệm cho
quan hệ người dùng--vai trò--quyền được sử dụng trong nguyên mẫu này. Nghiên cứu
này sử dụng RBAC lõi, chưa xét phân cấp vai trò hoặc ràng buộc phân tách nhiệm
vụ.

\subsection{Quản lý nhật ký bảo mật}
\label{subsec:log-management}

Nhật ký bảo mật cung cấp bằng chứng cho giám sát, điều tra và phản ứng. NIST SP
800-92 nhấn mạnh rằng tổ chức cần quy trình và hạ tầng để quản lý nhật ký bảo
mật máy tính một cách hiệu quả \citep{Kent2006SP80092}. Trong bài báo này, nhật
ký được xem như tập sự kiện có cấu trúc, có thể được phân tích, kiểm tra hợp lệ
và chuyển đổi thành phát hiện cũng như sự cố. Khía cạnh audit và accountability
cũng xuất hiện như một họ kiểm soát độc lập trong NIST SP 800-53
\citep{JTF2020SP80053}, nhưng nguyên mẫu chỉ minh họa việc phân tích log, không
nhằm chứng minh tuân thủ một bộ kiểm soát.

\subsection{Phát hiện theo luật và SQL injection}
\label{subsec:rule-detection}

SQL injection vẫn là một rủi ro phổ biến của ứng dụng web khi đầu vào người dùng
được đưa vào truy vấn cơ sở dữ liệu mà không tham số hóa đúng cách. OWASP khuyến
nghị sử dụng prepared statement, stored procedure, kiểm tra danh sách cho phép
và nguyên tắc đặc quyền tối thiểu như các biện pháp phòng vệ chính
\citep{OWASPSQLi}. Nguyên mẫu không thay thế các kiểm soát lập trình an toàn;
nó chỉ phát hiện một số chỉ báo SQL injection trong nhật ký nhằm phục vụ giám
sát và trình diễn.

\subsection{Giám sát dựa trên rủi ro}
\label{subsec:risk-monitoring}

Giám sát dựa trên rủi ro mở rộng quyết định nhị phân cho phép/từ chối bằng cách
xem xét ngữ cảnh như hành vi người dùng, IP nguồn, thời điểm và lịch sử sự kiện
gần đây. NIST SP 800-207 đặt trọng tâm bảo vệ vào người dùng, tài sản và tài
nguyên, thay vì giả định sự tin cậy dựa trên vị trí mạng \citep{Rose2020SP800207}.
NIST SP 800-63B cũng cho phép dùng chỉ báo gian lận, chẳng hạn IP hoặc vị trí
địa lý không mong đợi, làm tín hiệu cho kiểm soát xác thực bổ sung
\citep{NIST2025SP80063B}. Tổng quan hệ thống của Xiao và cộng sự cho thấy
RBAC có thể là nguồn ngữ cảnh trong thiết kế access control có xét rủi ro, dù
các tiếp cận tinh hạt thường thiên về ABAC \citep{Xiao2022ContextRisk}.

Nguyên mẫu trong bài báo không thay đổi quyết định cấp quyền tại thời gian thực.
Thay vào đó, nó dùng RBAC như một nguồn bằng chứng để hậu kiểm nhật ký và ưu
tiên hóa cảnh báo. Vì vậy, các tín hiệu IP mới, ngoài giờ, tài nguyên hiếm và
từ chối lặp lại được biểu diễn bởi heuristic theo luật có thể giải thích, không
phải một mô hình học máy hay một cơ chế access control thích nghi.

\subsection{Tương quan cảnh báo và phân loại sự cố}
\label{subsec:alert-correlation}

Một cảnh báo đơn lẻ hiếm khi mô tả đầy đủ một cuộc tấn công nhiều bước. Valdes
và Skinner nghiên cứu tương quan cảnh báo theo hướng xác suất
\citep{Valdes2001AlertCorrelation}; Zhou và cộng sự mô hình hóa quan hệ logic
giữa các alert để nhận diện alert cùng thuộc một intrusion
\citep{Zhou2007AlertCorrelation}. Trái với những mô hình correlation tổng quát
đó, RBAC Guard áp dụng quy tắc gom nhóm đơn giản theo cặp người dùng--IP và cửa
sổ thời gian. Lựa chọn này giảm phạm vi nhưng giúp evidence, IDs sự kiện và
điểm rủi ro của mỗi incident dễ kiểm tra trong môi trường học tập.

\subsection{Liên hệ với kịch bản tấn công}
\label{subsec:attack-scenarios}

Kịch bản đoán mật khẩu tương ứng với kỹ thuật Brute Force (T1110) trong ATT\&CK
\citep{MITRET1110}. Các chỉ báo SQL injection được dùng như dấu hiệu quan sát
trong request log; chúng không khẳng định khai thác thành công, nhưng phù hợp
với phạm vi giám sát một ứng dụng public-facing có thể bị khai thác
\citep{MITRET1190}. Các mapping này chỉ phục vụ mô tả kịch bản đánh giá, không
phải tuyên bố hệ thống bao phủ đầy đủ ATT\&CK.
